\section*{अध्याय \ref{ch:b2:sound-waves} --- तरलों में ध्वनि-तरंगें}

\begin{solution}{exo:b2:sound-waves:1}
$c = \sqrt{\gamma RT/M}$: \qty{331}{m/s}, \qty{343}{m/s}, \qty{299}{m/s}; हीलियम
\qty{1010}{m/s}; CO$_2$ \qty{268}{m/s}। स्वर-रज्जु उसी
आवृत्ति पर कंपते हैं, पर स्वर-नाल के अनुनाद (जो स्वर-गुण गढ़ते हैं)
$c$ के साथ बदलते हैं और तीन गुना ऊपर चले जाते हैं: अतः आवाज़ ऊँची
लगती है।
\end{solution}

\begin{solution}{exo:b2:sound-waves:2}
\qty{60}{dB}: $I = \qty{1e-6}{W/m^2}$, $p_m = \sqrt{2ZI} = \qty{29}{mPa}$, $v_m = p_m/Z =
\qty{7e-5}{m/s}$, $\xi_m = v_m/\omega = \qty{22}{nm}$। \qty{120}{dB}: \qty{1}{W/m^2},
\qty{29}{Pa}, \qty{7}{cm/s}, \qty{22}{\micro m}। दो: \qty{63}{dB}; सौ: \qty{80}{dB}।
\end{solution}

\begin{solution}{exo:b2:sound-waves:3}
$I = 1/4\pi r^2$: \qty{0.080}{W/m^2} (\qty{109}{dB}), \qty{89}{dB}, \qty{69}{dB}; और $I = 10^{-6}$
$r = \sqrt{1/4\pi10^{-6}} = \qty{280}{m}$ पर; \qty{30}{m} पर \qty{100}{dB} के लिए $4\pi \times
900 \times 10^{-2} = \qty{110}{W}$ ध्वनि चाहिए।
\end{solution}

\begin{solution}{exo:b2:sound-waves:4}
$4.5 \times 340 = \qty{1.5}{km}$; $340 \times 1.5 = \qty{510}{m}$; $1500 \times 1.2 = \qty{1.8}{km}$;
$1540 \times 65 \times 10^{-6} = \qty{10}{cm}$।
\end{solution}

\begin{solution}{exo:b2:sound-waves:5}
(क) $500 \times 340/300 = \qty{567}{Hz}$; $500 \times 340/380 = \qty{447}{Hz}$। (ख) $500 \times
380/340 = \qty{559}{Hz}$; दोनों: $500 \times 380/300 = \qty{633}{Hz}$, जबकि \qty{654}{Hz}
(अकेला स्रोत \qty{80}{m/s} पर) या \qty{618}{Hz} (अकेला प्रेक्षक): अर्थात् परिणाम हवा तय करती है,
सापेक्ष चाल नहीं। (ग) दीवार \qty{567}{Hz} पाती है
और उसे विराम में फिर उत्सर्जित करती है; और चालक \qty{40}{m/s} से पास आता है: $567 \times
380/340 = \qty{633}{Hz}$।
\end{solution}

\begin{solution}{exo:b2:sound-waves:6}
$v_m = p_m/Z$: \qty{2.4}{mm/s}, \qty{0.67}{\micro m/s}, \qty{25}{nm/s}; $\xi_m = v_m/\omega$:
\qty{390}{nm}, \qty{0.11}{nm}, \qty{4}{pm}; $I = p_m^2/2Z$: \qty{1.2e-3}{W/m^2} (\qty{91}{dB}),
\qty{3.3e-7}{W/m^2} (\qty{55}{dB}), \qty{1.3e-8}{W/m^2} (\qty{41}{dB})। कोई कड़ा माध्यम
दिए गए दाब पर मुश्किल से हिलता है और थोड़ी ही शक्ति ले जाता है।
\qty{1}{\micro Pa} के सापेक्ष \qty{1}{Pa} \qty{120}{dB} है।
\end{solution}

\begin{solution}{exo:b2:sound-waves:7}
(क) $\sqrt{1.013 \times 10^5/1.2} = \qty{290}{m/s}$; $\times\sqrt{1.4}$: \qty{343}{m/s}। (ख)
$\gamma = c^2M/RT$: हवा के लिए $1.40$, आर्गन के लिए $1.71$ ($5/3$, अंक-सन्निकटन के भीतर):
अर्थात् तीन स्थानांतरी स्वातंत्र्य-कोटियाँ ($\gamma = 1 + 2/3$) बनाम द्विपरमाणुक गैस की
पाँच ($1 + 2/5$)। (ग) $\sqrt{D/f} = \sqrt{2 \times 10^{-8}} = \qty{0.14}{mm}
\ll \lambda = \qty{34}{cm}$: अतः शिखरों और घाटियों के बीच कोई ऊष्मा नहीं बहती।
\end{solution}

\begin{solution}{exo:b2:sound-waves:8}
(क) खुली--खुली: $p = p_m\sin kx\cos\omega t$, जहाँ $kL = n\pi$; बंद--खुली:
$p = p_m\cos kx\cos\omega t$ (बंद सिरे पर वेग का निस्पंद, दाब का प्रस्पंद),
जहाँ $kL = (2n - 1)\pi/2$। (ख) $L = c/2f = \qty{39}{cm}$; $c/4f = \qty{19.5}{cm}$।
(ग) सभी संनादी; और केवल विषम संनादी (शहनाई-जैसी खोखली ध्वनि)।
(घ) प्रति खुले सिरे $0.6 \times 2 = \qty{1.2}{cm}$: $L_{\text{प्रभावी}} = \qty{41.4}{cm}$, $f =
\qty{414}{Hz}$, अर्थात् $6\%$ मंद --- नली को \qty{36.6}{cm} पर काटना पड़ेगा।
\end{solution}

\begin{solution}{exo:b2:sound-waves:9}
$\mathcal P = \qty{400}{W}$, $I = \mathcal P/2\pi r^2$। (क) \qty{6.4e-3}{W/m^2}, \qty{98}{dB},
$p_m = \qty{2.3}{Pa}$। (ख) $r = \sqrt{400/2\pi} = \qty{8}{m}$; और \qty{2.5}{km} पर \qty{70}{dB}।
(ग) $98 - 20\log(r/100) - (r - 100)/100 = 70$: $r \approx \qty{950}{m}$। (घ)
$\pi(10)^2/\pi(0.95)^2 \approx 110$ सायरन।
\end{solution}

\begin{solution}{exo:b2:sound-waves:10}
(क) $\rho_0\partial_tv = p_mk\sin kx\cos\omega t$: अतः $v = (p_m/Z)\sin kx\sin\omega t$। (ख)
$e_k = \tfrac12\chi_Sp_m^2\sin^2kx\sin^2\omega t$, $e_p = \tfrac12\chi_Sp_m^2\cos^2kx\cos^2\omega t$:
अर्थात् एक तब अधिकतम होता है जब दूसरा लुप्त हो, समय में भी और आकाश में भी। (ग) $E =
S\int_0^L(e_k + e_p)\dd x = \tfrac14\chi_Sp_m^2SL$, अचर। (घ) $I = pv =
(p_m^2/4Z)\sin2kx\sin2\omega t$: माध्य शून्य, पर हर आवर्तकाल में दो बार ऊर्जा
दाब-प्रस्पंदों से वेग-प्रस्पंदों तक और वापस बहती है।
\end{solution}

\begin{solution}{exo:b2:sound-waves:11}
(क) $t = 0$ पर उत्सर्जित अग्र की त्रिज्या $ct$ होती है जब स्रोत $vt$
दूर होता है: अतः स्पर्श-शंकु के लिए $\sin\theta = c/v$। (ख) $\theta = \qty{39}{\degree}$; और शंकु
वायुयान के नीचे ज़मीन तक तब पहुँचता है जब वह $h/\tan\theta = \qty{15}{km}$
आगे जा चुका हो: अर्थात् सिर के ऊपर से गुज़रने के $15000/480 = \qty{31}{s}$ बाद। (ग) पूरे
मार्ग के अनुदिश \qty{60}{km} चौड़ी पट्टी विस्फोट सुनती है। (घ) $\sin\theta
= 340/800$, $\theta = \qty{25}{\degree}$; पहले शंकु की तड़ाक आती है, और
नाल का धमाका, जो बंदूक से $c$ पर चलता है, बाद में।
\end{solution}

\begin{solution}{exo:b2:sound-waves:12}
(क) $f_1 = f(c + v\cos\theta)/c$। (ख) $f_2 = f_1c/(c - v\cos\theta) = f(c + v\cos\theta)
/(c - v\cos\theta) \approx f(1 + 2v\cos\theta/c)$। (ग) $2 \times 5 \times 10^6 \times 0.5 \times 0.5/1540
= \qty{1.6}{kHz}$; और \qty{90}{\degree} पर शून्य: अतः जाँच को तिरछा कीजिए। (घ) $\Delta v = 50 \times 1540/
(2 \times 5 \times 10^6 \times 0.5) = \qty{1.5}{cm/s}$। कान स्वर और स्वर-गुण दोनों का
तत्काल पीछा करता है: चिकनी सीटी \omterm{def:b2:viscous-flows:reynolds}{स्तरीय प्रवाह} बताती है, और फुफकार किसी संकरे स्थान के
पीछे का विक्षोभ।
\end{solution}

\begin{solution}{pb:b2:sound-waves:1}
\textbf{1.} $c = \sqrt{1.4 \times 8.31 \times 288/0.029} = \qty{340}{m/s}$; \qty{30}{\degreeCelsius}
पर \qty{349}{m/s}। \qty{2.9}{s} में \qty{1}{km}।

\textbf{2.} \qty{1.4}{km}। बिजली की नाली कई किलोमीटर लंबी होती है; उसके
भिन्न भागों की ध्वनि कई सेकंड तक आती रहती है, और साथ में बादलों तथा
ज़मीन से प्रतिध्वनियाँ भी।

\textbf{3.} $I = \qty{1}{W/m^2}$: $p_m = \sqrt{2ZI} = \qty{29}{Pa}$, $v_m = \qty{7}{cm/s}$।

\textbf{4.} $It = \qty{2}{J/m^2}$।

\textbf{5.} $c$ ऊपर की ओर बढ़ता है; और $\sin i/c$ के अचर रहने से $i$ बढ़ता है, अतः
किरण क्षैतिज की ओर और फिर नीचे मुड़ जाती है: इस तरह ध्वनि ज़मीन के
अनुदिश नलिकाबद्ध होकर दूर तक जाती है। पानी के ऊपर सबसे नीची हवा पानी से
ठंडी रहती है: वही बात।

\textbf{6.} किरणें ऊपर मुड़ जाती हैं और छाया-क्षेत्र छोड़ जाती हैं। वक्रता
त्रिज्या $R = c/(\dd c/\dd z) = \qty{50}{km}$: अतः \qty{3}{km} के बाद किरण
$3/50 = \qty{0.06}{rad} = \qty{3.4}{\degree}$ मुड़ चुकी होती है और \qty{90}{m} चढ़ चुकी होती है: अतः
श्रोता उसके नीचे रह जाता है।

\textbf{7.} अवशोषण $f^2$ की तरह बढ़ता है: कुछ किलोमीटर बाद तड़ाक की ऊँची
आवृत्तियाँ जा चुकी होती हैं और नीची गड़गड़ाहट बची रहती है।

\textbf{8.} $\mathcal P = \qty{2}{W}$; \qty{10}{m} पर $I = 2/2\pi r^2 = \qty{3.2e-3}{W/m^2}$:
अर्थात् \qty{95}{dB}।

\textbf{9.} $p_m = \sqrt{2 \times 410 \times 3.2 \times 10^{-3}} = \qty{1.6}{Pa}$; $v_m = \qty{3.9}{mm/s}$;
$\xi_m = v_m/2\pi f = \qty{6}{\micro m}$।

\textbf{10.} $+\qty{20}{dB}$: अर्थात् \qty{1}{m} पर \qty{115}{dB} --- जो मिनटों में
हानि पहुँचाता है। \qty{60}{dB}: $r = \sqrt{2/2\pi10^{-6}} = \qty{560}{m}$।

\textbf{11.} दस असंबद्ध: $+\qty{10}{dB}$, \qty{105}{dB}। कला में दो: दाब
दुगुना, $+\qty{6}{dB}$, \qty{101}{dB}।

\textbf{12.} संतुलन $\mathcal P = \alpha A\,ce/4$: $e = 8/(0.3 \times 2000 \times 340) =
\qty{3.9e-5}{J/m^3}$, $E = eV = \qty{0.4}{J}$। रुकने के बाद $\dd E/\dd t = -(\alpha Ac
/4V)E$: अतः $\tau = 4V/\alpha Ac = \qty{0.20}{s}$, और \qty{60}{dB} में $\tau\ln10^6 =
\qty{2.7}{s}$ लगते हैं (सैबीन का अनुरणन काल)।

\textbf{13.} $\lambda = c/f$: \qty{0.51}{mm}, \qty{0.31}{mm}, \qty{0.15}{mm}।

\textbf{14.} $2 \times 0.12/1540 = \qty{156}{\micro s}$; \qty{6.4}{kHz}; और \qty{31}{ms} में
$200$ रेखाएँ (अर्थात् तीस चित्र प्रति सेकंड)।

\textbf{15.} ऊतक--वायु: $r = -0.9995$, अर्थात् ऊर्जा का $99.9\%$ परावर्तित
--- इसीलिए जेल, जो हवा की झिल्ली हटा देता है। ऊतक--हड्डी: $r = 0.63$,
ऊर्जा का $40\%$; शेष हड्डी सोख लेती है, जो छाया
डालती है। वसा--पेशी: $r = 0.10$, ऊर्जा का $1\%$: अर्थात् वही मंद प्रतिध्वनियाँ जिनसे
चित्र बनता है।

\textbf{16.} परावर्तित अंश $1.1\%$ ($-\qty{20}{dB}$); \qty{5}{MHz} पर \qty{4}{cm}
की आवाजाही: $-\qty{10}{dB}$; कुल मिलाकर $-\qty{30}{dB}$, अर्थात् एक हज़ारवाँ भाग।

\textbf{17.} \qty{24}{cm} की आवाजाही: \qty{36}{dB}, \qty{60}{dB}, \qty{120}{dB}: अतः \qty{3}{MHz}
और \qty{5}{MHz} चलते हैं, \qty{10}{MHz} नहीं --- गहरे अंग कम आवृत्ति पर मोटे विभेदन के
साथ, और सतही अंग ऊँची आवृत्ति पर बारीकी से,
चित्रित किए जाते हैं।

\textbf{18.} पाँच तरंगदैर्घ्य, अर्थात् \qty{1.5}{mm} लंबा: अतः लगभग \qty{0.8}{mm}
(आधे स्पंद) से पास के दो अंतरापृष्ठ अतिव्यापी प्रतिध्वनियाँ देते हैं।

\textbf{19.} $2 \times 5 \times 10^6 \times 0.4 \times 0.5/1540 = \qty{1.3}{kHz}$: अर्थात् एक सीटी, जिसका
स्वर हर धड़कन के साथ चढ़ता और गिरता है।

\textbf{20.} $v = \qty{25}{m/s}$ पर: $700 \times 340/315 = \qty{756}{Hz}$, $700 \times 340/365 =
\qty{652}{Hz}$; अनुपात $1.16$, अर्थात् $12\log_21.16 = 2.6$ अर्ध-स्वर।

\textbf{21.} $700 \times 355/315 = \qty{789}{Hz}$।

\textbf{22.} निकटतम पहुँच पर, जब ऐंबुलेंस का वेग दृष्टि-रेखा के
लंबवत हो (ठीक-ठीक कहें तो, जब अभी आती ध्वनि उसी बिंदु पर
उत्सर्जित हुई हो)।

\textbf{23.} $\lambda' = (340 - 25)/700 = \qty{0.45}{m}$ (जबकि \qty{0.49}{m}); और
ध्वनि ऐंबुलेंस से आगे \qty{315}{m/s} से तथा पीछे \qty{365}{m/s} से
दूर जाती है।

\textbf{24.} $v = c\,\Delta f/2f = 3 \times 10^8 \times 3200/4.8 \times 10^{10} = \qty{20}{m/s} =
\qty{72}{km/h}$।

\textbf{25.} \qty{340}{m/s}, \qty{1540}{m/s}, \qty{25}{m/s}, \qty{20}{m/s}, \qty{0.4}{m/s},
\qty{3e8}{m/s}: $f' = f(c + v_o)/(c - v_s)$, एक या दो बार लगाया गया, और कम चाल पर
प्रकाश के लिए भी।
\end{solution}
